\chapter{Arsitektur Web Modern}
\label{ch:arsitektur}

\section*{Tujuan Pembelajaran}
\begin{itemize}
  \item Membedakan tempat HTML dirakit pada SPA, SSR, SSG, ISR, dan
        \textit{islands}, lalu menyebut akibatnya pada jumlah permintaan dan
        waktu sampai isi tampil.
  \item Mengukur kelima varian tersebut pada halaman yang sama, lalu membaca
        selisihnya dari byte yang tiba dan jumlah permintaan yang dibutuhkan.
  \item Memilih satu arsitektur untuk kebutuhan tertentu, lalu
        mempertanggungjawabkannya sebagai catatan keputusan berisi pilihan
        yang dipertimbangkan, yang diambil, dan konsekuensinya.
\end{itemize}

\section{Pendahuluan}

Bab 2 sampai Bab 7 memperbaiki satu aplikasi dari dalam: kuerinya, ingatannya,
kecepatannya, pengujiannya, dan pertahanannya. Bab ini mundur satu langkah dan
menanyakan hal yang lebih awal, yaitu di mana halaman sebaiknya dirakit.
Jawabannya menentukan banyak hal sesudahnya, termasuk berapa byte yang harus
diunduh pengunjung sebelum melihat isi.

Perbedaan antarpilihan itu sering dibahas sebagai selera, padahal sebagian
besarnya dapat dihitung. Satu halaman yang sama dapat disajikan dengan lima
cara, lalu keempat angkanya dibandingkan: byte yang tiba, apakah isinya sudah
ada di \textit{response} pertama, berapa permintaan lagi yang dibutuhkan, dan
berapa lama seluruhnya. Bab ini melakukan persis itu.

Demonya sengaja dibuat sekecil mungkin, yaitu satu file Python tanpa basis
data dan tanpa container, karena yang dibandingkan adalah tempat perakitannya,
bukan cara datanya diambil. Referensi utamanya adalah uraian
\textit{islands architecture} \cite{patterns-islands}, dokumentasi
\textit{server component} React \cite{react-server-components}, dokumentasi
ISR Next.js \cite{nextjs-isr}, dan catatan Fowler tentang memulai dari
\textit{monolith} \cite{fowler-monolith-first}.

\section{Dari Aplikasi Web Tradisional ke Single-Page Application}

Aplikasi web tradisional merakit HTML di server pada tiap permintaan. Browser
menerima halaman yang sudah jadi, menggambarnya, lalu mengulang seluruh
proses itu ketika pengunjung pindah halaman. Pola ini disebut
\textit{Server-Side Rendering} (SSR).

\textit{Single-Page Application} (SPA) membalik pembagiannya. Server mengirim
kerangka halaman yang hampir kosong beserta sebuah file JavaScript, lalu
browser yang meminta datanya dan menggambar isinya. Perpindahan halaman
sesudah itu tidak lagi meminta HTML baru, cukup data.

Analoginya seperti memesan makanan. SSR mengantar piring yang sudah terisi,
sedangkan SPA mengantar piring kosong beserta bahan dan kompornya, lalu
memasaknya di meja pemesan. Cara kedua membuat pesanan berikutnya jauh lebih
cepat, tetapi suapan pertama datang paling akhir.

Tabel \ref{tab:spa-ssr} memuat ukuran yang membedakan keduanya pada aplikasi
bab ini.

\begin{table}[!htb]
  \centering
  \caption{SSR lawan SPA pada halaman yang sama, hasil pengukuran nyata}
  \label{tab:spa-ssr}
  \begin{tabularx}{\textwidth}{Xrr}
    \hline
    \textbf{Ukuran} & \textbf{SSR} & \textbf{SPA} \\
    \hline
    Byte HTML yang pertama tiba & 511 & 229 \\
    \hline
    Byte seluruhnya sampai isi tampil & 511 & 1.551 \\
    \hline
    Permintaan sampai isi tampil & 1 & 3 \\
    \hline
    Isi katalog ada di HTML pertama & ya & tidak \\
    \hline
  \end{tabularx}
\end{table}

Baris terakhir yang paling menentukan. HTML pertama pada varian SPA memang
paling kecil, tetapi di dalamnya tidak ada satu judul buku pun. Isinya baru
ada sesudah browser mengunduh skripnya, menjalankannya, lalu meminta data
lewat permintaan ketiga.

\section[SSR, SSG, dan ISR]{Server-Side Rendering, Static Site
         Generation, dan Incremental Static Regeneration}

Ketiganya sama-sama mengirim HTML yang sudah berisi, dan hanya berbeda pada
kapan HTML itu dirakit. SSR merakitnya pada tiap permintaan.
\textit{Static Site Generation} (SSG) merakitnya sekali saat \textit{build},
lalu melayani file yang sama berulang kali. \textit{Incremental Static
Regeneration} (ISR) melayani file yang sudah jadi, tetapi merakitnya ulang
sesekali ketika umurnya lewat \cite{nextjs-isr}.

Gambar \ref{fig:tempat-rakit} merangkum ketiganya beserta dua varian lain
dalam satu garis.

\begin{figure}[!htb]
  \centering
  \begin{tikzpicture}[
    font=\scriptsize,
    kotak/.style={draw, rounded corners, minimum height=0.85cm,
                  minimum width=2.1cm, align=center, font=\scriptsize\bfseries},
    panah/.style={-{Stealth[length=2mm]}, thick}
  ]
    \node[kotak, fill=green!10] (build) at (0,0) {Saat build\\{\tiny\mdseries SSG}};
    \node[kotak, fill=orange!12] (sesekali) at (3.2,0) {Sesekali\\{\tiny\mdseries ISR}};
    \node[kotak, fill=blue!8] (tiap) at (6.4,0) {Tiap permintaan\\{\tiny\mdseries SSR, islands}};
    \node[kotak, fill=red!22] (browser) at (9.8,0) {Di browser\\{\tiny\mdseries SPA}};

    \draw[panah] (build) -- (sesekali);
    \draw[panah] (sesekali) -- (tiap);
    \draw[panah] (tiap) -- (browser);
    \node[font=\tiny, gray, anchor=north] at (4.9,-0.65)
      {makin ke kanan, makin segar isinya dan makin mahal tiap permintaannya};
  \end{tikzpicture}
  \caption{Kelima varian berbeda pada satu hal: kapan dan di mana HTML-nya
           dirakit}
  \label{fig:tempat-rakit}
\end{figure}

Tabel \ref{tab:varian} memuat pengukuran kelimanya pada halaman yang sama.
Perakitan halaman pada demo ini sengaja diberi jeda 50 milidetik, mewakili
halaman yang isinya mahal dihitung.

\begin{table}[!htb]
  \centering
  \caption{Lima varian pada halaman yang sama, hasil pengukuran nyata dengan
           \berkas{sources/chapter08/ukur-varian.py}}
  \label{tab:varian}
  \begin{tabularx}{\textwidth}{lrrlrr}
    \hline
    \textbf{Varian} & \textbf{HTML} & \textbf{Total} & \textbf{Isi di HTML} & \textbf{Minta} & \textbf{Total ms} \\
    \hline
    SSG & 511 B & 511 B & ya & 1 & 1,9 \\
    \hline
    ISR & 511 B & 511 B & ya & 1 & 52,9 \\
    \hline
    SSR & 511 B & 511 B & ya & 1 & 51,8 \\
    \hline
    Islands & 625 B & 1.322 B & ya & 2 & 3,1 \\
    \hline
    SPA & 229 B & 1.551 B & tidak & 3 & 1,7 \\
    \hline
  \end{tabularx}
\end{table}

Dua baris pada tabel itu perlu dibaca pelan. Baris ISR tercatat 52,9
milidetik karena pengukurannya kebetulan jatuh pada permintaan pertama
sesudah file lamanya kedaluwarsa, sehingga permintaan itulah yang menanggung
perakitan ulang. Baris SPA tercatat paling singkat justru karena yang diukur
hanya waktu jaringannya; pekerjaan menjalankan skrip dan menggambar isi
terjadi di browser dan tidak tertangkap alat ukur ini.

Listing \ref{lst:isr} memuat inti mekanismenya, yaitu delapan baris yang
memutuskan kapan file lama cukup dan kapan halamannya dirakit ulang.

\begin{lstlisting}[style=PythonStyle, caption={Inti Incremental Static Regeneration, tersedia di \berkas{sources/chapter08/aplikasi.py}}, label={lst:isr}]
def baca_atau_rakit_isr():
  """Membaca file ISR, dan merakitnya ulang bila umurnya sudah lewat.

  Inilah inti Incremental Static Regeneration: pengunjung tetap dilayani
  file yang sudah jadi, sedangkan perakitan ulang hanya terjadi sesekali.
  """
  if FILE_ISR.exists():
    umur = time.time() - FILE_ISR.stat().st_mtime
    if umur < UMUR_ISR_DETIK:
      return FILE_ISR.read_bytes(), "file lama dipakai"
  FILE_ISR.write_text(rakit_halaman_katalog(), encoding="utf-8")
  return FILE_ISR.read_bytes(), "dirakit ulang"
\end{lstlisting}

Tabel \ref{tab:isr} memperlihatkan perilaku ISR itu sendiri, diukur lima
permintaan berurutan dengan jeda sebelas detik sesudah permintaan ketiga.

\begin{table}[!htb]
  \centering
  \caption{Lima permintaan berurutan ke alamat ISR, umur file sepuluh detik,
           hasil pengukuran nyata}
  \label{tab:isr}
  \begin{tabularx}{\textwidth}{rXr}
    \hline
    \textbf{Permintaan} & \textbf{Keterangan dari header \texttt{X-ISR}} & \textbf{Waktu} \\
    \hline
    1 & dirakit ulang & 53,5 ms \\
    \hline
    2 & file lama dipakai & 1,9 ms \\
    \hline
    3 & file lama dipakai & 0,7 ms \\
    \hline
    4 & dirakit ulang, sesudah jeda sebelas detik & 51,5 ms \\
    \hline
    5 & file lama dipakai & 1,0 ms \\
    \hline
  \end{tabularx}
\end{table}

Pola itulah inti ISR. Biaya perakitan tidak hilang, melainkan ditanggung satu
pengunjung saja tiap satu periode, sedangkan sisanya dilayani secepat SSG.
Pertanyaan yang kemudian muncul: siapa pengunjung yang menanggungnya, dan
apakah itu dapat diterima. Bila tidak, perakitan ulang dipindahkan ke latar
belakang, sehingga pengunjung yang datang tepat sesudah kedaluwarsa tetap
menerima file lama.

\section{Islands Architecture}

\textit{Islands architecture} mengambil jalan tengah \cite{patterns-islands}.
Seluruh halaman dirakit di server seperti SSR, lalu hanya bagian yang
benar-benar perlu berinteraksi yang dihidupkan JavaScript. Bagian kecil itu
disebut pulau, dan sisanya tetap HTML biasa.

Pada demo bab ini pulaunya hanya satu, yaitu kotak penyaring judul. Daftar
bukunya sudah tergambar dari HTML yang dikirim server, sehingga pengunjung
yang tidak pernah mengetik di kotak itu tetap melihat katalog lengkap
walaupun skripnya gagal dimuat.

Tabel \ref{tab:islands} membandingkannya dengan SPA, yang menghidupkan
seluruh halaman.

\begin{table}[!htb]
  \centering
  \caption{Islands lawan SPA pada halaman yang sama, hasil pengukuran nyata}
  \label{tab:islands}
  \begin{tabularx}{\textwidth}{Xrr}
    \hline
    \textbf{Ukuran} & \textbf{Islands} & \textbf{SPA} \\
    \hline
    Byte HTML yang pertama tiba & 625 & 229 \\
    \hline
    Byte JavaScript yang diunduh & 697 & 938 \\
    \hline
    Permintaan sampai isi tampil & 2 & 3 \\
    \hline
    Isi tetap tampil bila skripnya gagal & ya & tidak \\
    \hline
  \end{tabularx}
\end{table}

Baris terakhir itu yang sering menentukan pilihan pada halaman publik.
Katalog, artikel, dan halaman produk tetap berguna tanpa JavaScript,
sedangkan papan kendali yang penuh interaksi memang tidak.

\section{Server Component dan Client Component}

Kerangka kerja modern membawa pembagian yang sama ke dalam satu bahasa.
\textit{Server component} dijalankan di server, hasilnya dikirim sebagai
HTML, dan kodenya tidak pernah ikut terunduh ke browser.
\textit{Client component} dikirim sebagai JavaScript dan dijalankan di
browser \cite{react-server-components}.

Tabel \ref{tab:komponen} merangkum pembagian tugasnya. Pembagian ini yang
membuat \textit{islands} dapat ditulis tanpa memisahkan dua proyek: komponen
yang perlu berinteraksi ditandai sebagai komponen klien, sisanya tetap di
server.

\begin{table}[!htb]
  \centering
  \caption{Pembagian tugas server component dan client component}
  \label{tab:komponen}
  \begin{tabularx}{\textwidth}{lXX}
    \hline
    \textbf{Hal} & \textbf{Server component} & \textbf{Client component} \\
    \hline
    Tempat berjalan & Server & Browser \\
    \hline
    Ikut diunduh & Tidak & Ya \\
    \hline
    Boleh menyentuh basis data & Ya & Tidak \\
    \hline
    Boleh memakai \textit{state} dan peristiwa & Tidak & Ya \\
    \hline
    Pantas dipakai untuk & Daftar, artikel, halaman detail & Tombol, formulir, penyaring \\
    \hline
  \end{tabularx}
\end{table}

\section{Edge dan Origin}

\textit{Origin} adalah server asal tempat aplikasi dan basis datanya berada.
\textit{Edge} adalah titik yang jauh lebih dekat ke pengunjung, dan
jumlahnya banyak. Bab 4 sudah memakai \textit{edge} dalam bentuk CDN untuk
file statis; bagian ini menanyakan apa lagi yang pantas dipindahkan ke sana.

Aturannya satu kalimat: yang pantas pindah ke \textit{edge} adalah pekerjaan
yang tidak butuh data tulis dan tidak butuh kedekatan dengan basis data.
Tabel \ref{tab:edge} memuat contohnya.

\begin{table}[!htb]
  \centering
  \caption{Pekerjaan yang pantas dan tidak pantas dipindahkan ke edge}
  \label{tab:edge}
  \begin{tabularx}{\textwidth}{Xl}
    \hline
    \textbf{Pekerjaan} & \textbf{Tempatnya} \\
    \hline
    Melayani file statis dan halaman hasil SSG & \textit{Edge} \\
    \hline
    Mengarahkan pengunjung menurut negara atau bahasa & \textit{Edge} \\
    \hline
    Memeriksa token sesi sebelum permintaan diteruskan & \textit{Edge} \\
    \hline
    Menulis pesanan dan memotong stok & \textit{Origin} \\
    \hline
    Menjalankan agregat yang menyentuh banyak tabel & \textit{Origin} \\
    \hline
  \end{tabularx}
\end{table}

Dua hal pada \textit{edge} sering tertukar. Pembagiannya begini. Menyimpan
salinan di \textit{edge} hanya memindahkan \textit{copy}, dan isinya dapat
\textit{stale} seperti dibahas Bab 4. Menjalankan kode di \textit{edge}
memindahkan perhitungannya, dan di sana jarak ke basis data justru bertambah
jauh. Karena itu kode yang banyak menyentuh basis data biasanya lebih cepat
tetap di \textit{origin}.

\section{Monolith, Modular Monolith, dan Backend for Frontend}

Tiga nama ini menjawab pertanyaan yang berbeda dari lima varian sebelumnya,
yaitu bagaimana kode sisi server dibagi. \textit{Monolith} menaruh semuanya
dalam satu aplikasi yang digelar sekali jalan. \textit{Modular monolith}
tetap satu aplikasi, tetapi dengan batas modul yang ditegakkan.
\textit{Backend for Frontend} (BFF) menambahkan satu lapisan tipis di depan,
khusus untuk satu jenis klien \cite{newman-bff}.

Gambar \ref{fig:bentuk-backend} memperlihatkan ketiganya.

\begin{figure}[!htb]
  \centering
  \begin{tikzpicture}[
    font=\scriptsize,
    kotak/.style={draw, rounded corners, minimum height=0.8cm,
                  minimum width=1.8cm, align=center, font=\scriptsize\bfseries},
    panah/.style={-{Stealth[length=2mm]}, thick}
  ]
    \node[kotak, fill=blue!8] (mono) at (0,0) {Monolith};
    \node[font=\tiny, anchor=north] at (0,-0.5) {satu aplikasi};

    \node[kotak, fill=orange!12, minimum width=2.6cm] (modular) at (3.6,0)
      {Modular monolith};
    \node[font=\tiny, anchor=north, align=center] at (3.6,-0.5)
      {satu aplikasi,\\batas modul ditegakkan};

    \node[kotak, fill=green!10] (bff1) at (7.6,0.55) {BFF web};
    \node[kotak, fill=green!10] (bff2) at (7.6,-0.55) {BFF mobile};
    \node[kotak, fill=blue!8] (inti) at (10.4,0) {Layanan inti};
    \draw[panah] (bff1) -- (inti);
    \draw[panah] (bff2) -- (inti);
    \node[font=\tiny, anchor=north, align=center] at (9,-1.1)
      {satu lapisan tipis\\untuk tiap jenis klien};
  \end{tikzpicture}
  \caption{Tiga cara membagi kode sisi server, dari yang paling sederhana ke
           yang paling banyak bagiannya}
  \label{fig:bentuk-backend}
\end{figure}

Urutan pada gambar tersebut juga urutan yang disarankan ketika ragu.
Memulai dari \textit{monolith} dan memecahnya ketika batasnya sudah
kelihatan lebih murah daripada memecah lebih dahulu lalu menyatukan kembali
\cite{fowler-monolith-first}. BFF dipakai ketika dua jenis klien benar-benar
menuntut bentuk data yang berbeda, misalnya aplikasi mobile yang butuh
jawaban jauh lebih ringkas daripada halaman web.

\section{Membaca dan Mempertanggungjawabkan Trade-off Arsitektur}

Keputusan arsitektur jarang punya pilihan yang benar sendiri. Yang ada adalah
pilihan yang konsekuensinya diketahui dan diterima. Karena itu keputusannya
ditulis, dan bentuk yang paling banyak dipakai adalah catatan keputusan
arsitektur, atau \textit{Architecture Decision Record} (ADR)
\cite{nygard-adr}.

Isinya pendek dan tetap: apa yang sedang diputuskan, pilihan apa saja yang
dipertimbangkan, mana yang diambil, dan apa akibatnya. Tabel \ref{tab:adr}
memperlihatkan satu catatan yang sudah terisi untuk halaman katalog bab ini.

\begin{table}[!htb]
  \centering
  \caption{Catatan keputusan untuk halaman katalog, memakai angka pada Tabel
           \ref{tab:varian}}
  \label{tab:adr}
  \begin{tabularx}{\textwidth}{lX}
    \hline
    \textbf{Bagian} & \textbf{Isinya} \\
    \hline
    Keputusan & Halaman katalog publik disajikan dengan ISR, umur sepuluh menit \\
    \hline
    Dipertimbangkan & SSR, SSG, dan ISR \\
    \hline
    Alasan & Isinya berubah beberapa kali sehari, sehingga SSG terlalu kaku, sedangkan SSR menanggung 51,8 ms pada tiap permintaan \\
    \hline
    Konsekuensi & Satu pengunjung tiap periode menanggung 52,9 ms, dan isi dapat \textit{stale} paling lama sepuluh menit \\
    \hline
    Yang membatalkan & Bila isi wajib segar seketika, atau bila perakitannya menjadi jauh lebih mahal \\
    \hline
  \end{tabularx}
\end{table}

Baris terakhir yang paling sering dilupakan. Menuliskan apa yang akan
membatalkan sebuah keputusan membuat keputusan itu dapat ditinjau ulang tanpa
perdebatan baru: cukup diperiksa apakah syarat pembatalnya sudah terjadi.

\section{Ringkasan}

Kelima varian arsitektur pada bab ini berbeda pada satu hal, yaitu kapan dan
di mana HTML dirakit. SSG merakitnya saat \textit{build}, ISR sesekali, SSR
pada tiap permintaan, \textit{islands} di server dengan sedikit bagian yang
dihidupkan di browser, dan SPA seluruhnya di browser.

Selisihnya dapat dihitung, bukan dikira-kira. Pada halaman yang sama, SSR
menanggung 51,8 milidetik tiap permintaan sedangkan SSG hanya 1,9 milidetik,
karena pekerjaan merakitnya sudah dibayar sekali saat \textit{build} dalam
0,3 milidetik. ISR mengambil keduanya: 53,5 milidetik pada permintaan
pertama sesudah kedaluwarsa, lalu 0,7 sampai 1,9 milidetik pada permintaan
berikutnya.

Ukuran yang paling menentukan bukan byte, melainkan ada atau tidaknya isi di
HTML yang pertama tiba. HTML varian SPA paling kecil, hanya 229 byte, tetapi
di dalamnya tidak ada satu judul pun, dan pengunjung menunggu tiga permintaan
sampai katalognya tampil. Varian \textit{islands} mengirim 625 byte yang
sudah berisi, lalu menambahkan satu skrip kecil hanya untuk kotak
penyaringnya.

Pembagian yang sama berlaku di sisi kode. \textit{Server component} tidak
ikut terunduh dan boleh menyentuh basis data, sedangkan \textit{client
component} diunduh dan boleh memakai \textit{state}. Di sisi penempatan,
\textit{edge} pantas menerima pekerjaan yang tidak menulis data dan tidak
butuh dekat dengan basis data, sedangkan sisanya tetap di \textit{origin}.

Keputusan arsitektur ditutup dengan catatan, bukan dengan kesepakatan lisan.
Catatan itu menyebut pilihan yang dipertimbangkan, yang diambil,
konsekuensinya, dan syarat yang akan membatalkannya. Dengan begitu
keputusannya dapat ditinjau ulang dengan memeriksa syarat, bukan dengan
mengulang perdebatan.

\section{Latihan}

Seluruh file yang dipakai pada latihan ini tersedia di
\berkas{sources/chapter08/}, dihitung relatif terhadap akar repositori.
Latihan ini hanya membutuhkan Python 3, tanpa Docker dan tanpa Node.js,
karena yang dibandingkan adalah tempat HTML dirakit.

Seluruh pengukuran diarahkan ke aplikasi yang berjalan di komputer sendiri,
dengan beban setara satu pengunjung membuka halaman. Hasil tiap mahasiswa
akan berbeda pada bagian waktu, dan yang dilaporkan adalah pola
perbandingannya beserta penjelasannya.

Penyiapannya dilakukan sekali, seperti pada Listing \ref{lst:c8-siapkan}.

\begin{lstlisting}[language=bash, caption={Membangun halaman statis lalu menyalakan server, memakai \berkas{sources/chapter08/bangun-statis.py}}, label={lst:c8-siapkan}]
cd sources/chapter08
python3 bangun-statis.py
python3 aplikasi.py &

# Keluarannya:
# Dibangun     : katalog-ssg.html
# Ukuran       : 511 byte
# Lama merakit : 0.3 ms
# Melayani di http://127.0.0.1:8040/ssr  (Ctrl+C untuk berhenti)
\end{lstlisting}

\subsection*{Latihan 1: Membandingkan Lima Varian}

Latihan ini membuktikan bahwa perbedaan kelima varian dapat dihitung, bukan
hanya diperdebatkan.

\begin{enumerate}
  \item Jalankan alat ukurnya seperti pada Listing \ref{lst:c8-lat1-jalan},
        lalu isi Tabel \ref{tab:c8-lat1-hasil}. Tabel
        \ref{tab:c8-lat1-contoh} memperlihatkan contoh yang sudah terisi.
  \item Buka kelima alamat di browser, lalu bandingkan apa yang tampil
        ketika JavaScript dimatikan lewat panel pengembang.
  \item Lihat sumber halaman tiap varian, lalu catat varian mana yang
        memuat judul buku di dalam HTML-nya.
  \item Jawab tiga pertanyaan berikut. Pertama, mengapa HTML varian SPA
        paling kecil tetapi totalnya paling besar? Kedua, mengapa alat ukur
        ini tidak dapat menangkap seluruh biaya varian SPA? Ketiga, varian
        mana yang paling tahan terhadap skrip yang gagal dimuat?
\end{enumerate}

\begin{lstlisting}[language=bash, caption={Mengukur kelima varian, tersedia di \berkas{sources/chapter08/ukur-varian.py}}, label={lst:c8-lat1-jalan}]
python3 ukur-varian.py

# Keluarannya:
#   Varian        HTML   Total  Isi di HTML  Minta  HTML ms  Total ms
#   ssg            511     511           ya      1      1.9       1.9
#   isr            511     511           ya      1     52.9      52.9
#   ssr            511     511           ya      1     51.8      51.8
#   islands        625    1322           ya      2      1.5       3.1
#   spa            229    1551        tidak      3      0.9       1.7
\end{lstlisting}

\begin{table}[!htb]
  \centering
  \caption{Contoh hasil Latihan 1 yang sudah terisi}
  \label{tab:c8-lat1-contoh}
  \begin{tabularx}{\textwidth}{lrrlr}
    \hline
    \textbf{Varian} & \textbf{HTML} & \textbf{Total} & \textbf{Isi di HTML} & \textbf{Minta} \\
    \hline
    SSG & 511 B & 511 B & ya & 1 \\
    \hline
    SPA & 229 B & 1.551 B & tidak & 3 \\
    \hline
  \end{tabularx}
\end{table}

\begin{table}[!htb]
  \centering
  \caption{Lembar isian hasil pengukuran Latihan 1}
  \label{tab:c8-lat1-hasil}
  \begin{tabularx}{\textwidth}{lrrlr}
    \hline
    \textbf{Varian} & \textbf{HTML} & \textbf{Total} & \textbf{Isi di HTML} & \textbf{Minta} \\
    \hline
    SSG & \dots & \dots & \dots & \dots \\
    \hline
    ISR & \dots & \dots & \dots & \dots \\
    \hline
    SSR & \dots & \dots & \dots & \dots \\
    \hline
    Islands & \dots & \dots & \dots & \dots \\
    \hline
    SPA & \dots & \dots & \dots & \dots \\
    \hline
  \end{tabularx}
\end{table}

\subsection*{Latihan 2: Membaca Perilaku ISR}

Latihan ini membuktikan bahwa biaya perakitan pada ISR tidak hilang,
melainkan ditanggung satu pengunjung tiap periode.

\begin{enumerate}
  \item Minta alamat ISR tiga kali berurutan, lalu catat header
        \texttt{X-ISR} dan waktunya, seperti pada Listing
        \ref{lst:c8-lat2-jalan}.
  \item Tunggu sebelas detik, lalu minta dua kali lagi dan catat kembali.
  \item Ubah \texttt{UMUR\_ISR\_DETIK} di
        \berkas{sources/chapter08/aplikasi.py} menjadi 60, nyalakan ulang
        aplikasinya, lalu ulangi pengamatannya.
  \item Isi Tabel \ref{tab:c8-lat2-hasil}. Tabel \ref{tab:c8-lat2-contoh}
        memperlihatkan contoh yang sudah terisi.
  \item Jawab dua pertanyaan berikut. Pertama, pengunjung mana yang
        menanggung biaya perakitan, dan apakah itu adil? Kedua, apa yang
        berubah bila umurnya dinaikkan menjadi satu jam?
\end{enumerate}

\begin{lstlisting}[language=bash, caption={Mengamati daur ISR, tersedia di \berkas{sources/chapter08/aplikasi.py}}, label={lst:c8-lat2-jalan}]
curl -s -D - -o /dev/null -w "%{time_total}\n" http://127.0.0.1:8040/isr \
  | grep -iE "^(x-isr|[0-9])"

# Keluarannya:
# X-ISR: dirakit ulang
# 0.053519
\end{lstlisting}

\begin{table}[!htb]
  \centering
  \caption{Contoh hasil Latihan 2 yang sudah terisi, umur file sepuluh detik}
  \label{tab:c8-lat2-contoh}
  \begin{tabularx}{\textwidth}{rXr}
    \hline
    \textbf{Permintaan} & \textbf{Keterangan \texttt{X-ISR}} & \textbf{Waktu} \\
    \hline
    1 & dirakit ulang & 53,5 ms \\
    \hline
    2 & file lama dipakai & 1,9 ms \\
    \hline
    4 & dirakit ulang, sesudah jeda & 51,5 ms \\
    \hline
  \end{tabularx}
\end{table}

\begin{table}[!htb]
  \centering
  \caption{Lembar isian hasil pengamatan Latihan 2}
  \label{tab:c8-lat2-hasil}
  \begin{tabularx}{\textwidth}{rXr}
    \hline
    \textbf{Permintaan} & \textbf{Keterangan \texttt{X-ISR}} & \textbf{Waktu} \\
    \hline
    1 & \dots & \dots \\
    \hline
    2 & \dots & \dots \\
    \hline
    3 & \dots & \dots \\
    \hline
    4, sesudah jeda sebelas detik & \dots & \dots \\
    \hline
    5 & \dots & \dots \\
    \hline
  \end{tabularx}
\end{table}

\subsection*{Latihan 3: Memindahkan Satu Bagian Menjadi Pulau}

Latihan ini membuktikan bahwa jumlah JavaScript yang diunduh ditentukan oleh
seberapa banyak bagian halaman yang dihidupkan.

\begin{enumerate}
  \item Baca \berkas{sources/chapter08/statis/pulau.js}, lalu catat
        ukurannya dan apa saja yang dikerjakannya.
  \item Bandingkan dengan \berkas{sources/chapter08/statis/spa.js}, lalu
        catat ukurannya dan apa yang akan hilang bila skrip itu gagal
        dimuat.
  \item Tambahkan satu pulau kedua pada halaman islands, misalnya tombol
        yang mengurutkan daftar menurut harga, lalu ukur ulang kedua varian.
  \item Isi Tabel \ref{tab:c8-lat3-hasil}. Tabel \ref{tab:c8-lat3-contoh}
        memperlihatkan contoh yang sudah terisi.
  \item Jawab dua pertanyaan berikut. Pertama, berapa byte JavaScript yang
        bertambah untuk tiap pulau baru? Kedua, pada titik mana menambah
        pulau menjadi tidak lagi lebih murah daripada memakai SPA?
\end{enumerate}

\begin{lstlisting}[language=bash, caption={Membandingkan ukuran skrip kedua varian, memakai \berkas{sources/chapter08/statis/}}, label={lst:c8-lat3-jalan}]
wc -c statis/pulau.js statis/spa.js

# Keluarannya:
#  697 statis/pulau.js
#  938 statis/spa.js
\end{lstlisting}

\begin{table}[!htb]
  \centering
  \caption{Contoh hasil Latihan 3 yang sudah terisi}
  \label{tab:c8-lat3-contoh}
  \begin{tabularx}{\textwidth}{Xrr}
    \hline
    \textbf{Varian} & \textbf{Byte skrip} & \textbf{Byte total} \\
    \hline
    Islands, satu pulau & 697 & 1.322 \\
    \hline
    SPA & 938 & 1.551 \\
    \hline
  \end{tabularx}
\end{table}

\begin{table}[!htb]
  \centering
  \caption{Lembar isian hasil pengukuran Latihan 3}
  \label{tab:c8-lat3-hasil}
  \begin{tabularx}{\textwidth}{Xrr}
    \hline
    \textbf{Varian} & \textbf{Byte skrip} & \textbf{Byte total} \\
    \hline
    Islands, satu pulau & \dots & \dots \\
    \hline
    Islands, dua pulau & \dots & \dots \\
    \hline
    SPA & \dots & \dots \\
    \hline
  \end{tabularx}
\end{table}

\subsection*{Latihan 4: Menulis Catatan Keputusan}

Latihan ini mengubah angka pada ketiga latihan sebelumnya menjadi satu
keputusan yang dapat dipertanggungjawabkan.

\begin{enumerate}
  \item Pilih satu kebutuhan nyata, misalnya halaman katalog publik, halaman
        rincian pesanan milik pengguna yang sudah masuk, atau papan kendali
        admin.
  \item Pilih tiga varian yang masuk akal bagi kebutuhan itu, lalu tulis
        angka pembandingnya dari Tabel \ref{tab:c8-lat1-hasil}.
  \item Tulis catatan keputusan berisi lima bagian seperti pada Tabel
        \ref{tab:c8-lat4-hasil}, yaitu keputusan, pilihan yang
        dipertimbangkan, alasan, konsekuensi, dan syarat yang
        membatalkannya. Tabel \ref{tab:c8-lat4-contoh} memperlihatkan
        contoh yang sudah terisi.
  \item Jawab dua pertanyaan berikut. Pertama, angka mana yang paling
        menentukan pilihan Anda, dan mengapa? Kedua, keadaan apa yang akan
        membuat keputusan itu perlu ditinjau ulang?
\end{enumerate}

\begin{table}[!htb]
  \centering
  \caption{Contoh catatan keputusan yang sudah terisi}
  \label{tab:c8-lat4-contoh}
  \begin{tabularx}{\textwidth}{lX}
    \hline
    \textbf{Bagian} & \textbf{Isinya} \\
    \hline
    Keputusan & Halaman katalog publik memakai ISR dengan umur sepuluh menit \\
    \hline
    Dipertimbangkan & SSR, SSG, ISR \\
    \hline
    Alasan & SSR menanggung 51,8 ms tiap permintaan, SSG terlalu kaku untuk isi yang berubah beberapa kali sehari \\
    \hline
    Konsekuensi & Satu pengunjung tiap periode menanggung 52,9 ms, isi dapat \textit{stale} sampai sepuluh menit \\
    \hline
    Yang membatalkan & Isi wajib segar seketika, atau perakitannya menjadi jauh lebih mahal \\
    \hline
  \end{tabularx}
\end{table}

\begin{table}[!htb]
  \centering
  \caption{Lembar isian catatan keputusan Latihan 4}
  \label{tab:c8-lat4-hasil}
  \begin{tabularx}{\textwidth}{lX}
    \hline
    \textbf{Bagian} & \textbf{Isinya} \\
    \hline
    Kebutuhan yang dipilih & \dots \\
    \hline
    Keputusan & \dots \\
    \hline
    Dipertimbangkan & \dots \\
    \hline
    Alasan, beserta angkanya & \dots \\
    \hline
    Konsekuensi & \dots \\
    \hline
    Yang membatalkan & \dots \\
    \hline
  \end{tabularx}
\end{table}

Setelah seluruh latihan selesai, server dihentikan dengan Ctrl+C. Tidak ada
container yang perlu dimatikan, karena bab ini memang tidak memakainya.
